\chapter{Alopecia Areata}

\section*{Definition}
Alopecia areata is an autoimmune, non-scarring hair loss disorder presenting as well-circumscribed patches of complete hair loss on the scalp, beard, or body. It affects 1--2\% of the population with a lifetime risk of 2\%, and can progress to total scalp (alopecia totalis) or whole-body (alopecia universalis) involvement.

\section*{What Happens in Your Body}
Alopecia areata is a T-cell-mediated autoimmune disease in which CD8+ NKG2D+ cytotoxic T lymphocytes target the hair follicle bulb (anagen stage), causing premature involution and hair shaft breakage. IFN-\ensuremath{\gamma} and IL-15 drive the immune attack, while JAK-STAT pathway signaling is central to the inflammatory cascade. The hair follicle remains viable (non-scarring), allowing potential for regrowth.

\section*{Causes}
\begin{itemize}
 \item \textbf{Genetic:} HLA class II alleles (DRB1*04, DQB1*03), ULBP3, and other loci; 10--20\% have family history
 \item \textbf{Autoimmune:} Associated with atopic disease (20--30\%), Hashimoto's thyroiditis, vitiligo, and type 1 diabetes
 \item \textbf{Triggers:} Severe emotional stress, physical stress (illness, injury), viral infection (suspected but unconfirmed)
 \item \textbf{Down syndrome:} Increased prevalence (5--10\% affected)
\end{itemize}

\section*{When to See a Doctor}

\begin{itemize}
 \item Acute onset of patchy hair loss, especially with smooth, non-scarred scalp surface
 \item Nail changes: pitting, trachyonychia (sandpaper nails), Beau's lines, red lunulae
 \item Rapid progression to extensive involvement, ophiasis pattern (band-like occipital)
 \item Psychological distress affecting quality of life
\end{itemize}

\section*{Related Symptoms}
\begin{itemize}
 \item Alopecia totalis, alopecia universalis, ophiasis, nail pitting, vitiligo, thyroid disease, atopic dermatitis
\end{itemize}
\textbf{Important:} The exclamation point hair (broken, tapering hairs) at the margin of active patches is pathognomonic on dermatoscopy and indicates active disease.

\section*{Self-Care / Home Management}
\begin{itemize}
 \item Use cosmetic covers: wigs, scarves, hairpieces, or eyebrow tattooing/microblading
 \item Apply sunscreen to exposed bald scalp areas to prevent sunburn
 \item Psychological support and counseling; connect with patient support groups (NAAF)
\end{itemize}

\section*{How Your Doctor Will Evaluate You}
\begin{itemize}
 \item Clinical diagnosis: Discrete, round patches of complete hair loss with smooth, non-scarred scalp; Pull test positive at active margins
 \item Dermatoscopy (trichoscopy): Exclamation point hairs, yellow dots (empty follicles), black dots (cadaverized hairs), broken hairs
 \item Scalp biopsy if diagnosis uncertain: Peribulbar lymphocytic infiltrate (\"swarm of bees\") in anagen follicles
 \item Screening for autoimmune disease: TSH, anti-TPO antibodies; atopy and thyroid disease evaluation
\end{itemize}

\section*{Treatment Options}
\begin{itemize}
 \item Limited patches (<50\% scalp): Intralesional triamcinolone acetonide (5--10 mg/mL per injection, 1 mL per site) every 4--6 weeks
 \item Children or needle-phobic: Topical corticosteroids (clobetasol 0.05\% foam/ointment) under occlusion
 \item Moderate--severe (>50\%): Topical immunotherapy with diphenylcyclopropenone (DPCP) or squaric acid dibutylester (SADBE)
 \item Systemic: Oral corticosteroids (limited bursts), methotrexate, cyclosporine; JAK inhibitors (tofacitinib, ruxolitinib, baricitinib) have shown efficacy and baricitinib is FDA-approved for severe alopecia areata
 \item Minoxidil 5\% topical BID as adjunctive therapy
\end{itemize}

\clearpage
